\subsection{情景结果与机制讨论}
在 2022 公报归一化口径下，我们先把基线情景看作恢复收益上传导的起点：江豚和长江鲟末值分别为 1.133 和 1.140。若只增强通道阻隔、暂不叠加额外污染，二者分别降至 0.901 和 0.856。再把长江鲟人工放流项置为 $R_A=0$，长江鲟末值降至 0.519，而江豚仅变为 1.145。污染情景下，江豚和长江鲟进一步降至 0.610 和 0.704。把这几组结果顺着看下来，通道约束会同时压低两类物种，而放流变化主要落在长江鲟上；污染则会在高阻隔基础上继续侵蚀基线收益。结果表聚焦江豚与长江鲟终态，$M$ 与 $V$ 则作为过程变量保留在后续机制解释中。
\begin{table}[htbp]
\centering
\caption{问题二的结果口径、相对变化与对照说明}
\label{tab:q2_result}
\begin{tabularx}{0.95\textwidth}{>{\raggedright\arraybackslash}p{2.2cm} >{\raggedright\arraybackslash}p{3.6cm} >{\raggedright\arraybackslash}p{2.8cm} >{\raggedright\arraybackslash}p{1.7cm} X}
\toprule
项目 & 数值/情景输出 & 相对基线变化 & 是否外部验证 & 解释用途 \\
\midrule
2022 公报基准 & 江豚 1249 头、长江鲟 438 尾 & -- & 是 & 作为问题二的初值与口径基准 \\
基线情景末值 & 江豚 1.133，长江鲟 1.140 & -- & 否 & 作为禁渔与食源修复下的基线投影 \\
$R_A=0$ 对照 & 江豚 1.145，长江鲟 0.519 & 江豚 +1.0\%，长江鲟 -54.5\% & 否 & 拆分自然恢复与人工放流补偿效应 \\
高阻隔对照 & 江豚 0.901，长江鲟 0.856 & 江豚 -20.5\%，长江鲟 -24.9\% & 否 & 仅提高通道阻隔，分离生境约束效应 \\
污染情景末值 & 江豚 0.610，长江鲟 0.704 & 江豚 -46.2\%，长江鲟 -38.2\% & 否 & 比较污染叠加阻隔对基线收益的侵蚀 \\
\bottomrule
\end{tabularx}
\end{table}
这几组对照摆在一起，机制差异就比较清楚了。无放流时，长江鲟大幅下降而江豚变化很小，说明人工补偿主要作用于鲟类而不是江豚；高阻隔与污染情景则同时压低两类物种，说明连通性和环境质量属于更上游的共性约束。作为定性验证，模型在 $t=3$（对应 2024 年）时长江鲟已呈现正增长趋势（$A(3)/A(0)>1$），与题目给出的“2024 年长江鲟自然繁殖群体重新出现产卵场”这一事实方向一致。

